\documentclass[10pt,twocolumn]{article}

\usepackage[margin=0.78in,columnsep=0.28in]{geometry}
\usepackage{times}
\usepackage{amsmath}
\usepackage{graphicx}
\usepackage{booktabs}
\usepackage{hyperref}
\usepackage[T1]{fontenc}
\usepackage{titlesec}
\usepackage{caption}
\usepackage{authblk}
\usepackage{enumitem}

\renewcommand{\dbltopfraction}{0.95}
\renewcommand{\dblfloatpagefraction}{0.8}
\setcounter{dbltopnumber}{3}
\titlespacing*{\section}{0pt}{8pt}{4pt}
\titlespacing*{\subsection}{0pt}{6pt}{3pt}
\captionsetup{font=small,labelfont=bf}

\title{\vspace{-1.2em}Does Fairness Improve RAG Performance, or Is It Diversity?\\
\large A Critical Reproduction and Extension of \emph{Towards Fair RAG}\vspace{-0.3em}}

\author[1]{Idan Peretz}
\author[1]{Avi Simkin}
\affil[1]{\normalsize \{308560150, 312485816\}}
\date{}

\begin{document}
\maketitle
\vspace{-2.2em}

\begin{abstract}
\noindent
Fairness-aware stochastic rerankers are usually expected to trade ranking quality for
more equitable item exposure, yet Eun Kim and Diaz (2025) report that RAG systems
built on such rerankers can maintain, and sometimes exceed, the generation utility of
deterministic ranking. We reproduce this setup on the LaMP benchmark and investigate
the mechanism behind it, contrasting a deterministic diversity-relevance re-ranker
(MMR) against pure stochastic sampling (Plackett-Luce) to separate the effect of
diversity from the effect of randomness itself. We find that relevance-aware
diversification is a real, though narrow and task-dependent, source of utility gain
over deterministic ranking, while blind randomization mostly is not --- evidence
against a simple story in which any diversification helps equally. Building on this,
we propose RandMMR, a hybrid reranker combining stochastic sampling with
relevance-aware diversification, and find it pushes fairness substantially further
than either ingredient alone while preserving utility --- an encouraging, though not
yet conclusively established, result. We also offer a critical reading of both the
original paper's evaluation protocol and of the LaMP benchmark itself as ground truth
for this kind of study, and lay out concrete directions for future work.
\end{abstract}

\vspace{0.3em}

\section{Introduction and Problem Formulation}
\label{sec:intro}

Fair-ranking research typically formalizes an unavoidable trade-off: enforcing more
equitable exposure among relevant items costs some ranking quality relative to a
purely relevance-ordered list \cite{castillo2018fairness,diaz2020evaluating}.
Eun Kim and Diaz \cite{eunkim2025fairrag} extend this literature to
retrieval-augmented generation (RAG), reporting that RAG systems using a stochastic,
fairness-aware retriever policy can maintain, and in some regimes even exceed, the
generation utility of a fully deterministic retriever, despite that trade-off. This
motivates three questions, the first two stated in our original project proposal:

\begin{itemize}[leftmargin=1.1em]
  \item \textbf{RQ1 (Reproduction).} Does the original paper's fairness-utility
  relationship hold under an independent re-implementation, and are there
  unmentioned limitations or sensitivities in how it is measured?
  \item \textbf{RQ2 (Mechanism).} If fairness interventions do not cost utility, is
  this because they increase the \emph{diversity} of retrieved items --- intuitively,
  a generator given a less redundant document set has more distinct evidence to draw
  on?
  \item \textbf{RQ3 (Improvement).} Given whatever RQ1--RQ2 reveal about \emph{why}
  fairness interventions do not cost utility, can that insight be used to design a
  reranking policy that pushes fairness (lower EE-D) further than the settings we
  study reach on their own, without giving utility back up in exchange?
\end{itemize}

Formally, following \cite{eunkim2025fairrag}, a RAG system is a tuple
$\langle R, G, k \rangle$ of a retriever $R$, generator $G$, and truncation depth
$k$. A deterministic $R$ returns one ranked list per query; a stochastic policy $S$
instead samples a ranking $\pi \sim S(\cdot \mid q)$, and system quality is reported
in expectation over $N$ sampled rankings: Expected Exposure Disparity/Relevance
(EE-D/EE-R, item-side ranking fairness and quality) and Expected Utility (EU,
downstream generation quality), each normalized to a theoretical or empirical bound
per query \cite{diaz2020evaluating,eunkim2025fairrag}. We add one further quantity
not defined in the source paper: \emph{intra-list diversity}
$\mathrm{ILD}(\pi) = 1 - \overline{\mathrm{Jaccard}}(\pi)$, the mean pairwise
token-set dissimilarity of the $k$ retrieved items, following the standard
similarity-vs-diversity framing of \cite{smyth2001similarity}. RQ2 asks whether EU is
better predicted by ILD than by the fairness dial itself.

\section{Related Work}
\label{sec:related}

\textbf{Fair and stochastic ranking.} Fairness in ranking is not simply fairness in
classification applied item-by-item: a ranked list's effect on an item is mediated by
\emph{exposure}, computed via the position-weighted formula defined by Diaz et al.\
\cite{diaz2020evaluating} rather than by an item's bare relevance. Item-side exposure
fairness is formalized via the Expected Exposure framework \cite{diaz2020evaluating}, which
measures a policy's exposure allocation against a fair target using a stochastic
ranking sampled by the Plackett-Luce (PL) model
\cite{plackett1975analysis,luce1959individual}; the temperature parameter $\alpha$
that we sweep here controls how close PL sampling is to the deterministic base
ranking. Castillo \cite{castillo2018fairness} surveys the fairness-transparency
trade-off in rankings that motivates this line of work. \emph{Towards Fair RAG}
\cite{eunkim2025fairrag} is the paper under reproduction: it applies this fairness
framework to RAG, measuring EE-D/EE-R/EU across LaMP tasks and several
retriever-generator pairs, and is the direct source of our experimental design,
metrics, and normalization convention (Appendix C of \cite{eunkim2025fairrag}). Its
utility metric, EU, is computed per sampled ranking as the downstream generator's
task-appropriate score (accuracy, MAE, or ROUGE-L, depending on the LaMP task) when
conditioned on that ranking's retrieved items, then normalized per query against an
oracle-ranking upper bound before being averaged across $N$ sampled rankings.

\textbf{Diversity in retrieval.} Maximal Marginal Relevance (MMR)
\cite{carbonell1998use} is the classic deterministic diversity-relevance trade-off
re-ranker, greedily selecting each next item to maximize $\lambda \cdot
\mathrm{relevance} - (1-\lambda)\cdot\max\text{-similarity-to-selected}$; we use it
as a second, non-stochastic diversification arm alongside PL. Intra-list diversity
via pairwise dissimilarity follows \cite{smyth2001similarity}.

\textbf{Personalization benchmark and models.} We use LaMP
\cite{salemi2023lamp} as our dataset (Section~\ref{sec:data}) and Flan-T5
\cite{chung2024scaling} as our generator family, matching \cite{eunkim2025fairrag}.

\section{Dataset}
\label{sec:data}

We use LaMP \cite{salemi2023lamp}, the same benchmark as \cite{eunkim2025fairrag},
for two reasons. First, and primarily, paper-faithfulness: LaMP is the only dataset
for which the original paper published per-task, per-generator utility-labeled
corpora (filtered to items with a measurable generator-specific utility delta), so
reproducing on it is the only way to compare our numbers directly against theirs
rather than against a different task distribution. Second, tractability: LaMP's
candidate pools are small enough that deriving actual relevance labels is
computationally feasible in the first place (Section~\ref{sec:threats}), which gives
us, for the same queries, both a retrieval-fairness signal (which documents are
labeled useful) and a generation-correctness signal --- a joint
retrieval-\emph{and}-generation labeling that,
to our knowledge, is not available together in other RAG-evaluation datasets, and is
exactly what studying fairness in RAG needs. LaMP additionally spans seven
personalization tasks with heterogeneous metrics (accuracy, MAE, ROUGE-L), letting us
check whether any finding is task-specific or general.

We verified our local corpus against the original paper's own published statistics
(Appendix D, Tables 4--5): per-task query count, average candidate-document count,
average useful-document count, and average useful-document percentage all match
to the paper's reported precision across all 7 tasks $\times$ 2 generators (28
cells), confirming our filtered dataset is the identical one used upstream, not a
re-derived approximation.

We run two generators (Flan-T5-Small, Flan-T5-Base \cite{chung2024scaling}) and two
retrievers (BM25, Contriever), against the paper's larger grid (adding SPLADE and
two larger generators) --- a scope reduction driven by available compute, not
methodology. Reranking settings are deterministic, a gold/oracle ranker, PL at
$\alpha \in \{1,2,4,8\}$, and MMR at $\lambda \in \{0.15,0.3,0.45,0.55,0.7,0.85\}$,
all truncated to $k=5$, matching \cite{eunkim2025fairrag}'s reported $k$.

\subsection{Limitations of the LaMP Ground Truth}
\label{sec:threats}

\begin{table*}[t]
\centering
\scriptsize
\caption{LaMP candidate-pool size and useful-document rate, per task and generator (from each task's relevance-mapping file).}
\label{tab:corpus}
\begin{tabular}{@{}lrrrrrrrrrrrrrr@{}}
\toprule
& \multicolumn{2}{c}{LaMP-1} & \multicolumn{2}{c}{LaMP-2} & \multicolumn{2}{c}{LaMP-3} & \multicolumn{2}{c}{LaMP-4} & \multicolumn{2}{c}{LaMP-5} & \multicolumn{2}{c}{LaMP-6} & \multicolumn{2}{c}{LaMP-7} \\
Generator & cand. & \%use & cand. & \%use & cand. & \%use & cand. & \%use & cand. & \%use & cand. & \%use & cand. & \%use \\
\midrule
Flan-T5-Small & 123.5 & 9.5 & 52.8 & 22.5 & 189.8 & 34.3 & 192.2 & 27.1 & 106.1 & 24.8 & 86.0 & 35.9 & 19.4 & 45.5 \\
Flan-T5-Base  & 102.9 & 22.5 & 45.6 & 29.9 & 185.3 & 41.2 & 187.0 & 31.6 & 105.6 & 25.7 & 86.2 & 38.7 & 21.7 & 33.0 \\
\bottomrule
\end{tabular}
\end{table*}

These reasons for choosing LaMP come with limitations worth stating up front, shared
with the source paper and not introduced by our reproduction, since they bound how
far any effect we measure in the results that follow (Section~\ref{sec:results}),
including RQ1's reproduction and RQ2's diversity test, should generalize.

\textbf{Relevance is a per-generator utility delta, not a fixed judgment.} Unlike
TREC-style collections, where relevance is a human judgment shared by every system, a
LaMP document is labeled relevant iff including it, versus a zero-context baseline,
\emph{improves that specific generator's} task score
($\delta = \text{augment\_score} - \text{baseline\_score}$; label $=1$ iff
$\delta{>}0$). Ground truth is therefore generator-specific: on the 21 LaMP-1 queries
shared by both generators' filtered sets, 22.8\% of document labels flip between
them, and the useful-document rate itself differs sharply (9.4\% vs.\ 20.2\%). Any
fairness-ranking result computed against these labels is a statement about one
model's learned usefulness, not a model-agnostic ground truth.

\textbf{The corpus is small, with an unusually high useful-document rate.} Table
\ref{tab:corpus} gives per-task, per-generator statistics: candidate pools of
19--192 documents/query (mean $\sim$110) and useful-document rates of
9.5\%--45.5\% (mean $\sim$29\%), against classical TREC collections' orders-of-
magnitude-larger pools and sub-1\% relevance rates. At $k{=}5$
\cite{eunkim2025fairrag} and a $\sim$30\% useful rate, even a fairness-motivated
rerank that displaces much of the top-$k$ still has decent odds of retaining several
useful documents, narrowing the deterministic-vs-stochastic utility gap the whole
fairness/utility trade-off is measured against.

\textbf{Our own generators are small; the source paper's are dated.} Flan-T5-Small
(77M) and Flan-T5-Base (250M) \cite{chung2024scaling}, the two we run, are one to two
orders of magnitude smaller than production RAG generators ($\geq$3B parameters).
The source paper's own larger generators are a separate limitation, not the same
one: its largest, Flan-UL2 ($\sim$20B), is not undersized so much as outdated ---
the same 2022-era family, not finetuned beyond LaMP's prompt format, and several
generations behind what is typically deployed today. Whether either paper's findings
persist at production scale is untested by both (see Future Work).

\textbf{The EE-D/EE-R trade-off is not a per-query guarantee.} The same small,
high-relevance corpus lets PL's aggregate trade-off fail query-by-query while holding
on average: an individual query's PL EE-R exceeds its own deterministic EE-R in
29.7\% of observations (37--38\% at $\alpha\in\{1,2,4\}$, still 8.3\% at the mildest
$\alpha{=}8$) --- though this mostly washes out once pooled to a macro mean (14 of
112 conditions, all at $\alpha{=}8$, $\Delta \leq 0.004$). A caution against
over-interpreting any single query's or run's numbers.

\section{Methodology Notes}
\label{sec:method}

We treat the original paper's open-sourced code as a direct reference throughout:
matching its library versions, reusing its utility-labeling technique
(Section~\ref{sec:threats}) and implementation details (prompt templates, scoring
functions) wherever our framework's structure allowed it directly, and otherwise
re-implementing to match its stated method as closely as we could. We report
divergences from the paper's own numbers (Section~\ref{sec:results}) as limitations
of exact reproduction under a reduced compute grid, not as bugs in either codebase,
since an extensive audit (prompt/scoring-function parity, generation parity against
the paper's own pipeline, model-version and seed invariance checks) found no
implementation-level explanation for them.

For the diversity-mediation test (RQ2), for each LaMP task we bin PL's and MMR's own
generated lists \emph{separately} into ILD$_{\text{norm}}$ quintiles and compare the
least-diversified quintile's (Q1) mean EU$_{\text{norm}}$ against the deterministic
baseline's, via a two-sided Welch's $t$-test, task-stratified with equal weight across
all four generator$\times$retriever cells. One MMR setting, $\lambda{=}1.0$, is
excluded: at $\lambda{=}1$ MMR's diversity term has zero weight, so its greedy
selection provably (and empirically, 100/100 on a real sample) reproduces the
deterministic ranking exactly --- it is not a diversification condition at all. A
second precision control matters at this scale: some (generator, ranker, $\alpha$)
conditions have been resampled more times ($N{\in}\{20,30\}$) than others ($N{=}10$),
and pooling each condition at its own maximum $N$ independently would silently give
some conditions tighter per-query utility estimates than the ones they are compared
against. Every PL-derived number in Section~\ref{sec:results} instead pools all
conditions inside one task at a single $N$ shared by every (generator, ranker,
$\alpha$) triple for that task --- the highest such $N$, or $N{=}10$ where no higher
value is universal.

\section{Results}
\label{sec:results}

\subsection{RQ1: Reproduction fidelity}
\label{sec:rq1}

Our reproduction of the paper's Figure 4 (EE-D$\leftrightarrow$EE-R$\leftrightarrow$EU
relationships) and Table 2 (disparity-interval utility deltas) recovers the same
qualitative trends the paper reports, once curve construction is matched exactly to
its stated method (binning every individual sampled list into disparity intervals
rather than averaging per $\alpha$-condition first, confirmed by pixel-level
inspection of the published figure): a fairness/ranking-quality trade-off that rises
then plateaus, a clean positive ranking-quality/utility relationship, and a similar
disparity-interval utility pattern to the paper's own Table 2. Absolute utility scale
is not directly comparable, since the paper's ceiling is a genuine stochastic oracle
(multiple sampled optimal permutations) while ours is a single deterministic oracle
draw.

Table~\ref{tab:baseline} (Appendix~\ref{app:baseline}) compares our reproduced
deterministic-baseline EU against the two tasks for which the paper publishes
baseline values directly. Flan-T5-Base is within 4\% on both. Flan-T5-Small
reproduces 22--30\% \emph{below} the paper on LaMP-1 --- traced to a genuine,
content-invariant model failure mode (it predicts the same class label regardless of
retrieved context on this task) --- and 38--42\% \emph{above} on LaMP-4, consistent
with our single-draw normalization ceiling being systematically looser than the
paper's true stochastic oracle. We treat both as reproduction-fidelity limitations
under a reduced compute grid, not as errors in either codebase.

\subsection{RQ2: Diversity as the mediating mechanism}
\label{sec:rq2}

\begin{figure}[t]
\centering
\includegraphics[width=\linewidth]{fig_diversity_lamp4_panel.png}
\caption{ILD-Jaccard diversity by rerank setting, LaMP-4 (representative; all seven
tasks show the same qualitative pattern, largest swing on LaMP-6 and smallest on
LaMP-7). Diamond = mean, orange line = median.}
\label{fig:diversity}
\end{figure}

Figure~\ref{fig:diversity} shows LaMP-4 as a representative example: PL and MMR both
raise list diversity as their dial is relaxed ($\alpha{:}8{\to}1$,
$\lambda{:}0.85{\to}0.15$). The other six tasks show the same qualitative pattern, but
the swing's size is task-dependent --- largest for LaMP-6 (mean ILD $0.64 \to 0.87$)
and smallest for LaMP-7 (everything within $0.92$--$0.96$), consistent with LaMP-7
having very few candidate documents per query ($\sim$20, Table~\ref{tab:corpus}) and
thus little room to diversify.

\begin{table*}[t]
\centering
\small
\caption{PL and MMR utility gain over deterministic, by each method's own
ILD$_{\text{norm}}$ quartile, LaMP-1 and LaMP-4 (the two tasks with a published
baseline). Bracketed values are the quartile's own ILD$_{\text{norm}}$ boundaries.
$^*p{<}0.05$, $^{**}p{<}0.01$ (Welch's $t$-test vs.\ deterministic).}
\label{tab:ild-quartiles}
\begin{tabular}{@{}llrllll@{}}
\toprule
Task & Method & Base.\ EU$_{\text{norm}}$ & Q1 (least diverse) & Q2 & Q3 & Q4 (most diverse) \\
\midrule
LaMP-1 & PL  & 0.626 & [0.71,0.95] $-$0.016 & [0.95,0.97] $-$0.105$^{**}$ & [0.97,0.98] $-$0.090$^{**}$ & [0.98,1.00] $-$0.129$^{**}$ \\
LaMP-1 & MMR & 0.626 & [0.82,0.96] $+$0.032 & [0.96,0.97] $+$0.004 & [0.97,0.99] $-$0.047 & [0.99,1.00] $-$0.044 \\
LaMP-4 & PL  & 0.304 & [0.45,0.97] $-$0.011 & [0.97,0.98] $-$0.048$^{**}$ & [0.98,0.99] $-$0.042$^{**}$ & [0.99,1.00] $-$0.045$^{**}$ \\
LaMP-4 & MMR & 0.304 & [0.45,0.96] $+$0.028$^{**}$ & [0.96,0.98] $+$0.006 & [0.98,0.99] $-$0.010 & [0.99,1.00] $-$0.001 \\
\bottomrule
\end{tabular}
\end{table*}

The key test is whether this diversity increase predicts \emph{higher} utility.
Table~\ref{tab:ild-quartiles} splits PL's and MMR's own lists into
ILD$_{\text{norm}}$ quartiles --- separately per method, unlike the original paper's
pooled EE-D-interval Table 2 --- and reports each quartile's actual boundary values
alongside its EU$_{\text{norm}}$ gain over deterministic, for the two tasks with a
published baseline (Table~\ref{tab:baseline}). PL is flat-to-negative at Q1 and
significantly negative at every coarser quartile, on both tasks. MMR is only ever
positive at Q1 (least diversified), significantly so for LaMP-4
($+0.028$, $p{<}0.01$); LaMP-1's Q1 is directionally positive ($+0.032$) but not
significant at this quartile split, though it \emph{is} significant at the finer
quintile split used in Figure~\ref{fig:mmr-vs-pl-q1} below --- a reminder that, on a
dataset and sample size this constrained, whether one specific cell clears
$p{<}0.05$ can depend on the exact binning choice, not just the underlying effect.

\begin{figure}[t]
\centering
\includegraphics[width=\linewidth]{fig_mmr_vs_pl_q1.pdf}
\caption{$\Delta$EU$_{\text{norm}}$ vs.\ deterministic at the Q1 (least-diversified)
bin, PL vs.\ MMR, per task. $^*p{<}0.05$, $^{**}p{<}0.01$.}
\label{fig:mmr-vs-pl-q1}
\end{figure}

Figure~\ref{fig:mmr-vs-pl-q1} breaks this down per task: MMR's Q1 delta vs.\
deterministic is positive in 6 of 7 tasks and significant in 4 (LaMP-3, LaMP-4,
LaMP-6 at $p{<}0.01$; LaMP-1 at $p{<}0.05$); the two non-significant positive tasks
are negligible (LaMP-5: $+0.0003$; LaMP-7: $+0.009$), and the one negative task is
also non-significant (LaMP-2: $-0.012$). PL's Q1 delta is positive in 4 of 7 tasks
and significant in just one (LaMP-3, $p{<}0.01$).
Where diversification helps at all, in other words, it is mostly MMR doing the work:
its relevance-aware selection retains more of what already made the original ranking
good, while PL's blind randomization mostly just adds noise.

This is a real but narrow finding, not a broad confirmation of the
diversity-mediation hypothesis in our original proposal: the benefit is concentrated
in a specific diversity range (mild, not maximal), is far stronger for MMR than PL,
and even for MMR is significant in barely more than half the tasks we test. It does
support a more specific claim: mild, relevance-aware diversification (MMR) is a real,
if task-dependent, source of utility gain over deterministic ranking, while PL's
fairness-motivated randomization, on its own, mostly is not.

\subsection{RQ3: RandMMR, a hybrid PL + randomized-MMR reranker}
\label{sec:rq3}

Motivated by RQ1--RQ2 --- fairness interventions do not reliably cost utility, and
mild MMR-style diversification is a real (if partial) reason why --- we propose
\textbf{RandMMR}: PL sampling of the top-ranked item, then a randomized-MMR pass
($\lambda \sim \mathrm{Unif}(\ell_{lo},\ell_{hi})$) over the remaining ranks. We
evaluate the RandMMR condition run across the most tasks and cells ($\alpha{=}4$,
$\lambda\sim\mathrm{Unif}(0.8,1.0)$, $N{=}20$: LaMP-$\{1,2,3,4,5,7\}$ $\times$ 4
generator/retriever cells, 24 completed conditions) against plain PL, choosing per
cell the PL $\alpha \in \{1,2,4,8\}$ whose own mean EE-D is closest to RandMMR's ---
biased toward an $\alpha$ with the same or slightly \emph{higher} disparity than
RandMMR's own, so PL never gets an unearned fairness advantage in the comparison ---
unless the single closest $\alpha$ is already within 0.02 EE-D, in which case that one
is used regardless of direction.

\begin{table*}[t]
\centering
\small
\caption{RandMMR ($\alpha{=}4$, $\lambda\sim U(0.8,1.0)$) vs.\ matched plain PL, LaMP-1
and LaMP-4 (macro means; PL $\alpha$ chosen per row, see text). EE-D is the paper's
own bounded metric; EU is our per-query empirical-max normalized EU$_{\text{norm}}$.
$^*p{<}0.05$.}
\label{tab:hybrid}
\begin{tabular}{@{}lllrrrrr@{}}
\toprule
Task & Generator & Ranker & PL $\alpha$ & PL EE-D/EU & RandMMR EE-D/EU & $\Delta$EU & Sig. \\
\midrule
LaMP-1 & T5-Base  & BM25       & 1 & 0.16/0.46 & 0.14/0.49 & $+$0.029 & \\
LaMP-1 & T5-Base  & Contriever & 1 & 0.16/0.47 & 0.13/0.49 & $+$0.027 & \\
LaMP-1 & T5-Small & BM25       & 2 & 0.11/0.11 & 0.12/0.11 & $+$0.000 & \\
LaMP-1 & T5-Small & Contriever & 1 & 0.15/0.10 & 0.12/0.13 & $+$0.029 & \\
LaMP-4 & T5-Base  & BM25       & 1 & 0.21/0.23 & 0.19/0.23 & $+$0.008 & \\
LaMP-4 & T5-Base  & Contriever & 1 & 0.21/0.23 & 0.18/0.24 & $+$0.013 & \\
LaMP-4 & T5-Small & BM25       & 2 & 0.20/0.18 & 0.20/0.23 & $+$0.045 & * \\
LaMP-4 & T5-Small & Contriever & 1 & 0.21/0.22 & 0.19/0.23 & $+$0.009 & \\
\bottomrule
\end{tabular}
\end{table*}

Table~\ref{tab:hybrid} details the two tasks with a published baseline; the chosen
$\alpha$ varies by generator/ranker pair (its own column, since it is not the same
across all four), and RandMMR's EU$_{\text{norm}}$ exceeds the matched PL's in 7 of
these 8 rows, significantly in one (T5-Small/BM25, LaMP-4, $p{<}0.001$). Across the
remaining 16 rows (LaMP-2, 3, 5, 7), it does so in 10 of 16 --- 17 of 24 completed
cells overall (71\%), with no cell showing a significant PL advantage. RandMMR's own
dial reaches further than the single $\alpha{=}4$ point shown here: at two additional
$\alpha$ values we ran for LaMP-1 ($8$ and $20$), EE-D drops as low as 0.11--0.14 for
the T5-Small cells --- below any workable plain-PL setting except $\alpha{=}1$, which
RandMMR matches or beats in utility at every $\alpha$ level tested.

We read this as encouraging, not conclusive: 24 completed cells is not the full
7-task$\times$4-cell grid, only one RandMMR configuration is complete enough for a
cross-task tally, and only one cell reaches significance individually. But the
direction is consistent (17 of 24, no significant losses): RandMMR is, at minimum,
not costing utility relative to comparably-fair plain PL, and appears to more often
help than not.

\section{Conclusions}
\label{sec:conclusion}

\textbf{RQ1.} The paper's central relationships (fairness vs.\ ranking quality;
ranking quality vs.\ utility) reproduce qualitatively once curve construction matches
the paper's method exactly, but absolute utility numbers diverge meaningfully
(Appendix~\ref{app:baseline}) despite an audit finding no implementation-level
explanation --- a limitation of exact reproduction under a reduced grid, not a
codebase error.

\textbf{RQ2.} We find a narrow, asymmetric form of support for the
diversity-mediation hypothesis in our proposal, not the broad, uniform version it was
originally framed as: diversification helps mostly through MMR's relevance-aware
selection, and only in about half the tasks we test, not through PL's
fairness-motivated randomization on its own (Section~\ref{sec:rq2}).

\textbf{RQ3.} RandMMR, our proposed hybrid, matches or exceeds matched plain PL's
utility in 17 of 24 tested cells while driving EE-D far below any workable plain-PL
setting except $\alpha{=}1$ (Section~\ref{sec:rq3}). Encouraging, not conclusive: the
grid is incomplete and only one cell is individually significant.

\section{Future Work}
\label{sec:future}

\begin{enumerate}[leftmargin=1.1em,itemsep=3pt]
  \item \textbf{Scale and realism.} Nearly every number in this paper comes
  from a setup chosen for what we could afford computationally, and as this is a course project, compromises were made:
  small sample counts, generators far below production scale, and a ground
  truth that depends on those same small generators' own judgments
  (Section~\ref{sec:threats}). Our own Flan-T5 models (77M/250M) are one to two
  orders of magnitude smaller than even the smallest generators commonly
  deployed in production; the source paper's largest, Flan-UL2 ($\sim$20B), is
  not undersized in the same way, but it is old --- the same 2022-era,
  non-RAG-finetuned family, several generations behind what is typically
  deployed today. A repeat at industry scale, with bigger sample counts and a
  dataset whose relevance labels do not depend on a single generator's own
  utility delta, would show how much of RQ1--RQ3 survives outside this regime.
  \item \textbf{Task-specific RAG fairness-diversity balancing.} An interesting course of research in our opinion would be to study how much diversity is beneficial for different tasks or queries and how it can be balanced with fairness constraints in an optimal way.
  \item \textbf{Other reranking designs.} RandMMR is one specific point in a
  larger design space --- PL sampling for exactly the top rank, MMR for
  everything after it. Comparing it against other splits of that same idea,
  and against diversity/fairness methods outside the PL/MMR family entirely
  (e.g.\ xQuAD-style sub-topic coverage), could be an interesting research direction and may possibly yield another diversity-fairness balancing algorithm.
\end{enumerate}

{\small
\begin{thebibliography}{99}

\bibitem{eunkim2025fairrag}
T.\ Eun Kim and F.\ Diaz.
\newblock Towards Fair RAG: On the Impact of Fair Ranking in Retrieval-Augmented
Generation.
\newblock \emph{ICTIR}, 2025. arXiv:2409.11598.

\bibitem{castillo2018fairness}
C.\ Castillo.
\newblock Fairness and Transparency in Ranking.
\newblock \emph{ACM SIGIR Forum}, 52(2), 2018.

\bibitem{diaz2020evaluating}
F.\ Diaz, B.\ Mitra, M.\ D.\ Ekstrand, A.\ J.\ Biega, and B.\ Carterette.
\newblock Evaluating Stochastic Rankings with Expected Exposure.
\newblock \emph{CIKM}, 2020.

\bibitem{salemi2023lamp}
A.\ Salemi, S.\ Mysore, M.\ Bendersky, and H.\ Zamani.
\newblock LaMP: When Large Language Models Meet Personalization.
\newblock arXiv:2304.11406, 2023.

\bibitem{chung2024scaling}
H.\ W.\ Chung et al.
\newblock Scaling Instruction-Finetuned Language Models.
\newblock \emph{JMLR}, 25, 2024.

\bibitem{carbonell1998use}
J.\ Carbonell and J.\ Goldstein.
\newblock The Use of MMR, Diversity-Based Reranking for Reordering Documents and
Producing Summaries.
\newblock \emph{SIGIR}, 1998.

\bibitem{smyth2001similarity}
B.\ Smyth and P.\ McClave.
\newblock Similarity vs.\ Diversity.
\newblock \emph{ICCBR}, 2001.

\bibitem{plackett1975analysis}
R.\ L.\ Plackett.
\newblock The Analysis of Permutations.
\newblock \emph{Journal of the Royal Statistical Society: Series C}, 24(2), 1975.

\bibitem{luce1959individual}
R.\ D.\ Luce.
\newblock \emph{Individual Choice Behavior: A Theoretical Analysis}.
\newblock Wiley, 1959.

\end{thebibliography}
}

\clearpage
\appendix
\section{Deterministic-Baseline EU Detail}
\label{app:baseline}

\begin{table}[h]
\centering
\small
\caption{Deterministic-baseline EU: paper vs.\ our reproduction (Section~\ref{sec:rq1}).}
\label{tab:baseline}
\begin{tabular}{@{}llrrr@{}}
\toprule
Task & Model & Paper & Ours & $\Delta$\% \\
\midrule
LaMP-1 & BM25+Small       & 0.308 & 0.240 & $-$22\% \\
LaMP-1 & Contr.+Small     & 0.286 & 0.200 & $-$30\% \\
LaMP-1 & BM25+Base        & 0.670 & 0.695 & $+$4\% \\
LaMP-1 & Contr.+Base      & 0.637 & 0.641 & $+$1\% \\
LaMP-4 & BM25+Small       & 0.217 & 0.299 & $+$38\% \\
LaMP-4 & Contr.+Small     & 0.254 & 0.360 & $+$42\% \\
LaMP-4 & BM25+Base        & 0.223 & 0.304 & $+$36\% \\
LaMP-4 & Contr.+Base      & 0.268 & 0.361 & $+$35\% \\
\bottomrule
\end{tabular}
\end{table}

\end{document}
